\subsection{Пространственные методы анализа EEG}

Многоканальный характер электроэнцефалографических данных позволяет рассматривать сигнал не только во временной и частотной, но и в пространственной области. Пространственные методы направлены на построение линейных преобразований, усиливающих информативные компоненты сигнала и подавляющих шум и нерелевантную активность. В задачах декодирования EEG такие методы часто играют ключевую роль, особенно при низком отношении сигнал/шум.

\subsubsection*{Common Spatial Patterns}

Метод общих пространственных паттернов (Common Spatial Patterns, CSP) является одним из наиболее известных подходов к пространственной фильтрации в задачах бинарной классификации EEG~\cite{10.1109/86.895947}. Основная идея CSP заключается в поиске линейных комбинаций каналов, максимизирующих дисперсию сигнала для одного класса и минимизирующих её для другого. 

Формально задача сводится к решению обобщённой задачи на собственные значения для ковариационных матриц двух классов. Полученные фильтры проецируют исходный сигнал в подпространство, где различия между состояниями становятся более выраженными. Метод продемонстрировал высокую эффективность в задачах моторного воображения (motor imagery), в том числе в онлайн-сценариях.

Однако CSP чувствителен к артефактам и требует корректной оценки ковариационных матриц. Кроме того, классическая версия ориентирована на бинарную классификацию и плохо масштабируется на многоклассовые задачи без дополнительных модификаций.

% ---- КАРТИНКА 2 ----
% Иллюстрация принципа CSP:
% исходные каналы → линейная комбинация → усиление различий дисперсий двух классов

\subsubsection*{Алгоритм xDAWN}

Для задач, основанных на вызванных потенциалах (event-related potentials, ERP), предложен алгоритм xDAWN~\cite{10.1109/TBME.2009.2012869}. В отличие от CSP, ориентированного на дисперсионные различия, xDAWN направлен на усиление слабых ERP-компонентов, таких как P300, за счёт оптимального выбора пространственных фильтров.

Метод оценивает подпространство, в котором отношение сигнал/шум для целевого компонента максимально, что позволяет повысить качество классификации при ограниченном числе усреднений. Пространственная фильтрация в xDAWN снижает размерность данных и усиливает информативные компоненты перед подачей в классификатор.

\subsubsection*{Ковариационные представления и риманова геометрия}

Альтернативный подход к пространственному анализу основан на использовании ковариационных матриц многоканального EEG как основных признаков. Ковариационная матрица симметрична и положительно определена (symmetric positive definite, SPD), что означает, что пространство таких матриц обладает неевклидовой геометрией.

Риманова геометрия (Riemannian geometry) предоставляет естественный аппарат для сравнения SPD-матриц с использованием геодезических расстояний на соответствующем многообразии~\cite{10.1080/2326263X.2017.1297192}. Вместо работы в обычном евклидовом пространстве элементы данных рассматриваются как точки на римановом многообразии, что позволяет корректно учитывать геометрические свойства ковариаций.

Данный подход демонстрирует высокую устойчивость при малом числе обучающих примеров и хорошую переносимость между сессиями и субъектами. Современные работы рассматривают римановский фреймворк как перспективную основу для построения устойчивых BCI-систем нового поколения~\cite{10.48550/arXiv.1310.8115}. 

Использование ковариационных представлений особенно оправдано в условиях ограниченного объёма данных, когда глубокие нейронные сети склонны к переобучению. Кроме того, такой подход естественным образом учитывает пространственную структуру сигнала без явного построения фильтров вручную.

% ---- КАРТИНКА 3 ----
% Схема: ковариационная матрица как точка на римановом многообразии SPD
% Евклидово расстояние vs риманово геодезическое расстояние

\subsection{Спектральные и временно-частотные признаки}

Помимо пространственной структуры, информативность EEG-сигнала часто проявляется в частотной области. Осцилляторная активность отражает изменения синхронизации нейронных ансамблей и может быть чувствительна к когнитивному состоянию.

\subsubsection*{Band power и ERD/ERS}

Изменения мощности в определённых частотных диапазонах описываются в терминах событийно-связанной десинхронизации и синхронизации (event-related desynchronization / synchronization, ERD/ERS)~\cite{10.1016/S1388-2457-99-00141-8}. В отличие от ERP, которые фазово синхронизированы с событием, ERD/ERS отражают изменения амплитуды продолжающейся ритмической активности.

Частотные диапазоны альфа и бета традиционно связываются с процессами внимания, моторного контроля и памяти, что делает анализ band power обоснованным в задачах визуального воображения и реконструкции образов.

% ---- КАРТИНКА 4 ----
% Иллюстрация ERD/ERS:
% временная ось + изменение мощности в альфа-диапазоне относительно baseline

\subsubsection*{Оценка спектра мощности}

Для оценки спектральных характеристик применяются методы оценки спектра мощности, такие как метод Уэлча (Welch’s method) и вейвлет-преобразование (wavelet transform). Метод Уэлча основан на усреднении спектров по перекрывающимся окнам и обеспечивает устойчивую оценку мощности в заданном диапазоне частот. Вейвлет-анализ позволяет получать временно-частотные карты (time-frequency representations), отражающие динамику спектра во времени.

\subsection{Временные представления и агрегация}

EEG-сигнал обладает выраженной временной динамикой, поэтому использование скользящих окон (sliding windows) и анализ отдельных проб (single-trial analysis) являются важными инструментами повышения устойчивости декодирования. Применение временной агрегации и постобработки предсказаний, например через анализ наиболее длинной последовательной повторяемости метки, может повысить стабильность классификации~\cite{10.1109/JSEN.2023.3270281}.

\subsection*{Выводы по разделу}

Современные методы извлечения признаков EEG включают пространственные фильтры (CSP, xDAWN), ковариационные представления с использованием римановой геометрии, а также спектральные и временно-частотные признаки. Пространственные и ковариационные подходы позволяют эффективно учитывать многоканальную структуру сигнала, тогда как спектральный анализ отражает функциональную организацию нейронной активности. В задачах реконструкции визуальных стимулов комбинирование пространственно-спектральных признаков представляется наиболее обоснованным подходом.